\section{自然演绎与逆否式}


\begin{quote}
来源说明：本人选课较晚，本节主要依据作业与参考教材整理，非课堂记录。
\end{quote}


\subsection{记号与两种“反证”动作}


采用课程约定：$\neg\varphi$ 是 $\varphi\to\bot$ 的缩写。

\begin{itemize}
\item \textbf{$\to I$}：在临时假设 $[\varphi]^i$ 下推出 $\psi$，解除该假设得到 $\varphi\to\psi$。若推出的是 $\bot$，得到 $\neg\varphi$。
\item \textbf{经典 RAA}：在临时假设 $[\neg\varphi]^i$ 下推出 $\bot$，解除它得到 $\varphi$。
\item \textbf{$\to E$}：由 $\varphi\to\psi$ 和 $\varphi$ 得 $\psi$。
\end{itemize}


所以“假设 $\varphi$ 导出矛盾”得到的是 $\neg\varphi$，而“假设 $\neg\varphi$ 导出矛盾”才使用 RAA 得到 $\varphi$。

\subsection{逆否式}


证明

\[
(\neg Q\to\neg P)\leftrightarrow(P\to Q).
\]

正向推导：

\begin{verbatim}
1.  ¬Q→¬P        AS
2.    P           AS
3.      ¬Q        AS
4.      ¬P        →E 1,3
5.      ⊥         →E 2,4
6.    Q           RAA 3
7.  P→Q           →I 2-6
8.  (¬Q→¬P)→(P→Q) →I 1-7
\end{verbatim}


反向推导：

\begin{verbatim}
1.  P→Q           AS
2.    P           AS
3.    Q           →E 1,2
4.      ¬Q        AS
5.      ⊥         →E 3,4
6.    ¬P          →I 2-5
7.  ¬Q→¬P         →I 4-6
8.  (P→Q)→(¬Q→¬P) →I 1-7
\end{verbatim}


编号只表示要解除的假设。第一段的 RAA 解除第 3 行的 $[\neg Q]$；第二段的第 6 行是由 $[P]$ 导出矛盾，所以必须写 $\to I$。

\subsection{Fitch 列表与 Reiteration}


典型证明：

\begin{verbatim}
1 | ¬¬P        PR
2 |   ¬P       AS
3 |   ⊥        →E 2,1
4 | P          RAA 2
\end{verbatim}


缩进表示临时假设的作用域。已经写过且仍在作用域中的公式可以用 \texttt{R, 行号} 重复引用；已关闭子证明中的行不能再引用。
